\subsection{Equal-budget update paths: represented gradients and black-box search}
\label{sec:native-method-paths-v10}

V9 shows that implementing a backward pass does not establish a useful search
direction. We therefore define a prospective v10 comparison of two update
policies, rather than retrospectively replacing its no-op by the better known
endpoint. Both vary the same charter evidence weight $\theta\in[.2,.8]$ in the
original five-author, two-year, seed-42 society with fixed Qwen3:8B weights and
the same real-paper corpus. The native integer review parser, scientific
activities and selection thresholds remain unchanged. Both optimize exactly
$L=3.25A-Q+.05R$ from Eq.~\eqref{eq:native-qwen-v9-loss}, without a hidden
merit, fairness or promotion veto.

\paragraph{Posthoc diagnosis, not new native evidence.}
Recomputing the frozen v9 representation at different temperatures exposes a
value-fit/gradient trade-off (Table~\ref{tab:native-v10-diagnosis}). The sharpest
relaxation fits calibration losses closely but is locally flat at the center;
coarser relaxations even change the gradient's sign. Only one of 18 fitted raw
review responses has a substantial coordinate slope. Its influence passes
through both individual integer-score and admission relaxations; at $.01$
both are nearly saturated. Cutting represented memory backward changes the
center derivative only at numerical precision. This diagnosis uses no new
model, container or native-world calls and does not repair v9's failed unused
prediction. A larger represented gradient need not improve native loss.

\begin{table}[t]
\caption{Offline ablations of the same frozen v9 calibration representation
at $\theta=.5$. These are posthoc model evaluations, not four native
experiments or evidence for choosing a retrospectively successful update.}
\label{tab:native-v10-diagnosis}
\centering
\small
\begin{tabular}{rrr}
\toprule
Temperature $\tau$ & Center $\partial\widehat L/\partial\theta$ & Calibration-loss RMSE \\
\midrule
1    & $-.24092$ & $.52717$ \\
.2   & $-.09426$ & $.32974$ \\
.05  & $+.01156$ & $.01522$ \\
.01  & $1.15\!\times\!10^{-15}$ & $1.16\!\times\!10^{-5}$ \\
\bottomrule
\end{tabular}
\end{table}

\paragraph{Same information and loss, different proposal rules.}
Each arm receives only the three historical v9 calibration endpoints
$\theta=(.2,.5,.8)$ with native losses $(.9,.9,.8167)$, and is charged three
historical evaluations. Their underlying calls are counted once, not twice.
Both start their \emph{current} coordinate at $.5$ while already having the
historical \emph{incumbent} $.8$. The incumbent is the lowest observed native
loss, with earliest-observation ties; it is not substituted for the current
iterate. Neither optimizer receives old unused endpoints or the other arm's
new outcomes. All these historical data are development knowledge, not a new
held-out sample.

At gradient step $k=0,1,2$, refit the original affine raw-response operator on
the shared three endpoints plus only its own earlier completed endpoints.
Use the fixed temperature schedule $\tau_k=(1,.2,.05)$, not RMSE-based selection
among temperatures. Backpropagate at the current coordinate and propose
\begin{equation}
 \theta_{k+1}=\operatorname{clip}_{[.2,.8]}
 \bigl(\theta_k+\operatorname{clip}_{[-.1,.1]}(-.5\widehat g_k)\bigr).
 \label{eq:native-v10-step}
\end{equation}
The raw gradient is not sign-normalized or replaced when flat. A no-op still
requires its allocated fresh evaluation, and even a worsening result advances
the current iterate. The smoothed representation never changes the hard native
optimization objective.

The black-box arm uses no derivative or response representation. Sort distinct
observed coordinates, find the intervals immediately adjacent to its incumbent,
and evaluate the midpoint of the widest interval; equal widths choose the
lower interval. Add only its own new observation and repeat. This is a simple
local interval-refinement baseline, not an optimal global-search method.

\paragraph{Budget and interpretation.}
The precommitted order is G0, B0, G1, B1, G2, B2: three fresh native attempts
per arm. Proposal, eligible-history and operator hashes are persisted before
execution. Technical failure has null loss and stops the campaign without a
replacement. Each complete original-native forward shares the verified
53-reply historical prefix, then makes affected year-2 review and later calls
freshly. Logical responses, fresh calls and historical-origin calls are
separate ledgers. An explicit source-equivalence bridge checks original
scientific module bytes and compatible controls without rewriting v9 records.
Paths must distinguish current, proposed and incumbent systems, and improvement
over the initial center from improvement over the historical best. Equal native
budgets do not imply equal response costs. One coupled society, three updates
and potentially stochastic same-coordinate replies cannot establish optimizer
superiority, significance, causal institutional benefit or an exact
native/Qwen Jacobian.
